% Generated draft: requires mathematical and source review.
\documentclass[11pt]{article}
\usepackage[margin=25mm]{geometry}
\usepackage{fontspec}
\setmainfont{Latin Modern Roman}
\usepackage{amsmath,amssymb,mathtools,unicode-math}
\setmathfont{Latin Modern Math}
\usepackage{enumitem,needspace,xurl,hyperref}
\setlength{\parindent}{0pt}
\setlength{\parskip}{5pt}
\setlength{\emergencystretch}{3em}
\newcommand{\N}{\mathbb{N}}
\newcommand{\Z}{\mathbb{Z}}
\newcommand{\Q}{\mathbb{Q}}
\newcommand{\R}{\mathbb{R}}
\newcommand{\C}{\mathbb{C}}
\newcommand{\F}{\mathbb{F}}
\newcommand{\A}{\mathbb{A}}
\newcommand{\PP}{\mathbb{P}}
\newcommand{\E}{\mathbb{E}}
\newcommand{\eps}{\varepsilon}
\newcommand{\dd}{\,\mathrm{d}}
\begin{document}
\section*{UnsolvedMath ID: 10800001}
\textbf{Source problem code:} catalogue entry
\subsection{Definitions and mathematical statement}
For an integer $\mu\geq1$, let $\operatorname{UConf}_{\mu}(\C)$ be the space of unordered sets of $\mu$ distinct complex numbers. With a base configuration $q$, its fundamental group is the braid group $B_\mu=\pi_1(\operatorname{UConf}_{\mu}(\C),q)$.

In the usual singularity-theoretic setting, the Lyashko--Looijenga map sends a deformation parameter to the unordered critical values of the deformed function, with $\mu$ the relevant Milnor number. Write
\[
 \operatorname{LL}:E=\C^\mu\setminus\Sigma
 \longrightarrow \operatorname{UConf}_{\mu}(\C)
\]
for the restriction intended in the source. The singularity, deformation, precise excluded locus $\Sigma$, and local versus global parameter domain are not supplied in this excerpt.

Vassiliev's Problem 1A asks for explicit braid-group obstructions to lifting by this map, and specifically to characterize which braids cannot be lifted into $E$. For a chosen $a\in E$ with $\operatorname{LL}(a)=q$, a path-lifting formulation asks for a criterion on $\beta\in B_\mu$ for the existence of paths $b:[0,1]\to\operatorname{UConf}_{\mu}(\C)$ and $\gamma:[0,1]\to E$ such that
\[
 [b]=\beta,\qquad b(0)=b(1)=q,\qquad
 \gamma(0)=a,\qquad \operatorname{LL}\circ\gamma=b.
\]
A lift need not be a closed path. Requiring $\gamma(1)=a$ instead would be a different question about closed lifts. The source must determine whether liftability concerns a fixed initial parameter, some initial parameter, or specified representatives of a braid class; these distinctions matter when the map is not a covering.
\subsection{Short English statement}
Find explicit braid-group criteria obstructing lifts through a Lyashko--Looijenga map. The singularity, excluded parameter locus, and basepoint and path-lifting conventions must be specified.
\subsection{Sources}
\begin{itemize}
\item Local registry: \url{data/unsolvedmath/problems.json}, permanent UnsolvedMath ID $10800001$, exact source code \texttt{catalogue entry}, attributed to Victor Vassiliev.
\item Vassiliev, \emph{A Few Problems on Monodromy and Discriminants}, Problem 1A: \url{https://amj.math.stonybrook.edu/html-articles/Files-2015-2024/15-11/}. Supplied access date: 2026-07-29; extraction: html-structure. The article contains eighteen explicitly labeled problem environments.
\item The registry reports the Lyashko--Looijenga covering theorem for simple singularities and nonliftability phenomena for nonsimple ones, attributing the latter to nonproperness. No separate primary references or exact hypotheses are supplied for those assertions.
\end{itemize}
\end{document}
